\section{Conclusion}
\label{sec:Con}

This study evaluated a minimal runtime verification pipeline for automotive ECU telemetry replay, using LARVA-driven monitoring over four properties (\texttt{voltageOutOfRange}, \texttt{coolantTooHigh}, \texttt{throttleSpike}, \texttt{loadMapMismatch}). The empirical objective was to quantify monitoring overhead and detection behaviour under controlled baseline and abnormal traces.

The findings indicate that: (i) injected abnormal classes are detected consistently when RV is enabled and are not reported when RV is disabled; (ii) monitoring overhead is low for nominal traces (3.41\%) but substantial in violation-heavy scenarios (692.36\%); and (iii) threshold tuning reduces false positives without collapsing abnormal detection performance. The controlled three-variant comparison showed that abnormal overhead is mainly driven by generated bad-state diagnostics: 3927.05\% (original generated) decreased to 683.79\% (logging disabled) and then to 9.39\% after removing stacktrace-based bad-state string construction, with CPU-time measurements showing the same trend. The hypothesis is therefore \textbf{partially supported}: useful safety/cybersecurity coverage is achievable with bounded overhead for nominal workloads, and worst-case overhead can be reduced substantially through targeted monitor-code optimisation.

Methodological limitations should be acknowledged. First, evaluation is trace-driven and server-side, not hardware-in-the-loop on target automotive microcontrollers. The assumed deployment is that IoT/ECU records are forwarded to a server for verification, so measured overhead is server-path overhead rather than on-device monitor overhead. Second, the selected property set is intentionally small and does not cover all dynamics of production powertrain/braking stacks. Third, the abnormal dataset is synthetic (though controlled), and therefore may not represent full real-world fault diversity.

Future work should focus on four directions. (1) Move from server replay to ECU-representative hardware and measure CPU/memory budgets directly. (2) Introduce adaptive/specification-level optimisations (e.g., event filtering, staged monitors, lower-cost violation handling) to reduce violation-path latency. (3) Expand validation with recorded CAN/FlexRay traces and hardware-in-the-loop campaigns. (4) Integrate trust primitives (e.g., attestation and hardware trust anchors) with monitor lifecycle management to align runtime verification with broader automotive cybersecurity assurance frameworks~\cite{ekert2021cybersecurity,plappert2023hta,lorych2024hta}.